\section{Examples and limits}\label{sec:example}

\paragraph{A connected graph.} Take $G_1$ of \cref{sec:counting} ($N=3$, edges $v_1v_2$, $v_2v_3$), whose $A$ and $A^2$ are printed there. Then
\[
B=A+A^2=\begin{pmatrix}1&1&1\\1&2&1\\1&1&1\end{pmatrix}.
\]
Every off-diagonal entry of $B$ is positive, so $G_1$ is connected by \cref{thm:sum}. The pairs $v_1v_2$ and $v_2v_3$ get their positive entry from $A$, which counts the one-step walks $v_1v_2$ and $v_2v_3$. The pair $v_1v_3$ gets $B_{13}=0+1=1$ from $(A^2)_{13}=1$, which counts the walk $v_1v_2v_3$; this walk is also a path.

\paragraph{A disconnected graph.} Let $G_2$ have $N=3$, the edge $v_1v_2$, and the isolated vertex $v_3$. Then
\[
A=\begin{pmatrix}0&1&0\\1&0&0\\0&0&0\end{pmatrix},\quad
A^2=\begin{pmatrix}1&0&0\\0&1&0\\0&0&0\end{pmatrix},\quad
B=\begin{pmatrix}1&1&0\\1&1&0\\0&0&0\end{pmatrix}.
\]
The entries $B_{13}$, $B_{23}$, $B_{31}$, $B_{32}$ are $0$, so $G_2$ is not connected by \cref{thm:sum}. This agrees with the graph: $v_3$ has no neighbour, so no walk leaves it.

\paragraph{The sum must reach $A^{N-1}$.} Let $G_3$ have $N=4$ and edges $v_1v_2$, $v_2v_3$, $v_3v_4$. Then
\[
A=\begin{pmatrix}0&1&0&0\\1&0&1&0\\0&1&0&1\\0&0&1&0\end{pmatrix},\;
A^2=\begin{pmatrix}1&0&1&0\\0&2&0&1\\1&0&2&0\\0&1&0&1\end{pmatrix},\;
A^3=\begin{pmatrix}0&2&0&1\\2&0&3&0\\0&3&0&2\\1&0&2&0\end{pmatrix}.
\]
The entry $(1,4)$ is $0$ in $A$ and in $A^2$ (the only neighbour of $v_1$ is $v_2$, which is not adjacent to $v_4$, so no walk of length $2$ ends at $v_4$) and is $1$ in $A^3$, counting the walk $v_1v_2v_3v_4$. A sum stopped at $A^2$ would have $(1,4)$ entry $0$ and so wrongly report $G_3$ as disconnected; the path $v_1v_2v_3v_4$ has length $3=N-1$ (\cref{lem:reach}). The full sum is
\[
B=A+A^2+A^3=\begin{pmatrix}1&3&1&1\\3&2&4&1\\1&4&2&3\\1&1&3&1\end{pmatrix},
\]
whose off-diagonal entries are all positive (the diagonal entries are $B_{ii}$), so $G_3$ is connected by \cref{thm:sum}.

\paragraph{Limits.} The test takes only $A$ as input, and it has the following limits. It reads only whether entries are zero: a positive entry of $A^k$ means a walk of length exactly $k$ exists. Hence (\cref{lem:cut}) a path of length at most $k$ exists, not necessarily one of length exactly $k$ (in $G_1$, whose $A^3$ is printed in \cref{sec:counting}, $(A^3)_{12}=2$, yet the only path from $v_1$ to $v_2$ has length $1$). The diagonal of $B$ is left out on purpose: a connected graph can have zeros there, for example $N=2$ with one edge, where $B=A$ has a zero diagonal, so requiring a positive diagonal would reject a connected graph. The test is proved here only for finite simple undirected graphs with $N\ge2$.
